\emph{Setting.} By Lemma \claimref{singularity.pseudo-inverse-identities-hyperbolic-atom}, $v = \tfrac12(e_{0} + e_{1})$ and $v^{*} = \tfrac12(e_{0} - e_{1})$ satisfy the canonical anticommutation relations $v^{2} = (v^{*})^{2} = 0$, $vv^{*} + v^{*}v = 1$ in $\Cl(1,1) = \mathrm{alg}(e_{0}, e_{1}) \subset \Cl(1,3)$, and by Theorem \claimref{singularity.modular-structure-completed-pair}(1) and (3), $\sigma = v^{*}v - vv^{*} = e_{0}e_{1}$, $\sigma^{2} = 1$, and $\tfrac12(1 \pm \sigma)$ are $v^{*}v$ and $vv^{*}$. Directly, $\sigma v = e_{0}e_{1}\tfrac12(e_{0} + e_{1}) = \tfrac12(-e_{1} - e_{0}) = -v$ and $v\sigma = v$, so $\sigma$ anticommutes with $v$, and likewise $\sigma v^{*} = v^{*} = -v^{*}\sigma$.

\emph{(i).} $J^{2} = e_{2}e_{3}e_{2}e_{3} = -e_{2}^{2}e_{3}^{2} = -1$. Each of $e_{0}$, $e_{1}$ anticommutes with $e_{2}$ and with $e_{3}$ (orthogonal vectors), hence commutes with their product $J$; so $J$ commutes with every element of $\Cl(1,1)$. $J\sigma = e_{2}e_{3}e_{0}e_{1}$: moving $e_{0}$ and then $e_{1}$ to the front past two factors each takes four transpositions, so $J\sigma = e_{0}e_{1}e_{2}e_{3} = I_{4}$, and $\sigma J = I_{4}$ as $J$ and $\sigma$ commute; with $\sigma^{2} = 1$, $J = I_{4}\sigma$. The frame $*$ of Theorem \claimref{tower.structure-tower-c-norms} on the real algebra is the antiautomorphism fixing the frame generator of square $+1$ and negating those of square $-1$; $x \mapsto e_{0}\rev{x}e_{0}$ is an antiautomorphism (reversion is one, conjugation by $e_{0}$ an automorphism) with $e_{0} \mapsto e_{0}$ and $e_{k} \mapsto e_{0}e_{k}e_{0} = -e_{k}$, so it is that $*$. Then $J^{*} = e_{3}^{*}e_{2}^{*} = e_{3}e_{2} = -J$.

\emph{(ii).} $\Cl(1,1)$ has basis $1, e_{0}, e_{1}, \sigma$, and the eight products of these with $1$ and $J$ are distinct basis blades of $\Cl(1,3)$, so $\Aa = \Cl(1,1) \oplus \Cl(1,1)J$ has dimension $8$; it is closed under products because $J$ commutes with $\Cl(1,1)$ and $J^{2} = -1$. An element $c = c_{1} + c_{2}J$ ($c_{1}, c_{2} \in \Cl(1,1)$) is central in $\Aa$ exactly when $c_{1}$ and $c_{2}$ are central in $\Cl(1,1)$, whose centre is $\R$ (even dimension), so the centre of $\Aa$ is $\mathrm{span}\{1, J\}$, a field isomorphic to $\C$ as in any real algebra $\mathrm{span}\{1, u\}$ with $u^{2} = -1$. Over that field $\Aa$ is spanned by $1, e_{0}, e_{1}, \sigma$, and $e_{0} = v + v^{*}$, $e_{1} = v - v^{*}$, $\sigma = 2v^{*}v - 1$, so also by $1, v, v^{*}, v^{*}v$. Over any field $F$ the algebra presented by $v, v^{*}$ with the canonical anticommutation relations maps onto $\Mat_{2}(F)$ by its Fock representation ($v$ the matrix unit $E_{12}$, $v^{*}$ the matrix unit $E_{21}$), and it is spanned by $1, v, v^{*}, v^{*}v$ (every word in $v, v^{*}$ reduces to these by $v^{2} = 0$ and $vv^{*} = 1 - v^{*}v$); $\Aa$ is a quotient of it of dimension $4$ over $F = \mathrm{span}\{1, J\}$, equal to $\dim \Mat_{2}(F)$, so $\Aa \cong \Mat_{2}(F)$, the CAR algebra of one mode over $F$. The $*$ preserves $\Aa$ ($e_{0}$, $e_{1}$, $J$ are mapped into it), $v^{*} = \tfrac12(e_{0} - e_{1})$ is the image of $v$, $\sigma^{*} = e_{1}^{*}e_{0}^{*} = -e_{1}e_{0} = \sigma$, and $J^{*} = -J$ with $J$ central gives $(zx)^{*} = \bar{z}x^{*}$ for $z = a + bJ$, $\bar{z} = a - bJ$: conjugate-linear. The script checks the dimension, the closure, the centre by solving the linear system for the commutant of $e_{0}, e_{1}, J$ in $\Aa$, the spanning set $\{1, v, v^{*}, v^{*}v\}$ over the field by rank, and the antiautomorphism property of $*$ on all $256$ pairs of blades.

\emph{(iii).} $I_{4} = J\sigma$ lies in $\Aa$. Being the pseudoscalar of the even-dimensional $\Cl(1,3)$, it anticommutes with every vector, in particular with $e_{0}$, $e_{1}$ and their combinations $v$, $v^{*}$, so it is not central in $\Aa$. For vectors $a, b$, $(I_{4}a)(I_{4}b) = -I_{4}^{2}ab = ab$, while $J$ commutes with $a$ and $b$ of the radial plane and $(Ja)(Jb) = J^{2}ab = -ab$. $I_{4}P_{\pm} = J\sigma\tfrac12(1 \pm \sigma) = J\tfrac12(\sigma \pm 1) = \pm JP_{\pm}$. The script checks these on $e_{0}$, $e_{1}$, $v$, $v^{*}$ and $3e_{0} - 2e_{1}$, all pairs.

\emph{(iv).} $JK = \tfrac12J\sigma = \tfrac12I_{4}$ and $(2JK)^{2} = I_{4}^{2} = -1$. For any element $A$ with $A^{2} = -1$ the exponential series collapses as in Theorem \claimref{std-algebra.bivector-exponential} ($A^{2k} = (-1)^{k}$, $A^{2k+1} = (-1)^{k}A$; there for a bivector, the argument uses only the scalar square), so $e^{\theta A} = \cos\theta + A\sin\theta$, and with $A = 2JK$, $\theta = t/2$: $e^{JtK} = \cos\tfrac{t}{2} + 2JK\sin\tfrac{t}{2}$. Since $J$ commutes with $v$ and $\sigma v = -v$, $JKv = -\tfrac12Jv$ and $vJK = \tfrac12Jv$, so $JKv - vJK = -Jv$; likewise $+Jv^{*}$ for $v^{*}$. As $JK$ anticommutes with $v$, $e^{JtK}v = ve^{-JtK}$, so $\tau_{t}(v) = e^{JtK}ve^{-JtK} = e^{2JtK}v = e^{I_{4}t}v$, and $I_{4}v = J\sigma v = -Jv$ with $J$ commuting with $v$ gives $e^{I_{4}t}v = (\cos t - J\sin t)v = e^{-Jt}v$, the exponential $e^{-Jt} = \cos t - J\sin t$ being taken in the field $\mathrm{span}\{1, J\}$. For $v^{*}$, $I_{4}v^{*} = J\sigma v^{*} = Jv^{*}$ and $\tau_{t}(v^{*}) = e^{+Jt}v^{*}$. The $*$ is an antiautomorphism with $(JK)^{*} = K^{*}J^{*} = -KJ = -JK$, so $(e^{JtK})^{*} = e^{-JtK} = (e^{JtK})^{-1}$: $e^{JtK}$ is unitary and $\tau_{t}$ is a $*$-automorphism, $\tau_{t}\tau_{s} = \tau_{t+s}$, and it maps $\Aa$ into itself because $e^{JtK} \in \Aa$. $Jv = \tfrac12(e_{2}e_{3}e_{0} + e_{2}e_{3}e_{1})$ is a trivector, so for $\sin t \neq 0$ the image of a nonzero vector $a = \alpha e_{0} + \beta e_{1}$ of the plane, $e^{I_{4}t}a = a\cos t + I_{4}a\sin t$, has a nonzero trivector part. The script checks the commutators, $JK = I_{4}/2$ and $(2JK)^{2} = -1$, and at three rational points $(c, s) = ((1 - u^{2})/(1 + u^{2}), 2u/(1 + u^{2}))$ of the circle with $U = c + 2sJK$ (that is $c = \cos\tfrac{t}{2}$, $s = \sin\tfrac{t}{2}$): $UU^{-1} = 1$ with $U^{-1} = U^{*} = c - 2sJK$, $UvU^{-1} = (C - SJ)v$ and $Uv^{*}U^{-1} = (C + SJ)v^{*}$ with $C = c^{2} - s^{2}$, $S = 2cs$, against the wrong sign $(C + SJ)v$, and the addition law for two such points.

\emph{(v).} For $z = t + J\theta \in \mathrm{span}\{1, J\}$, $JzK$ commutes with $K$ and with $J$, and the computation of (iv) goes through with $z$ in place of $t$ because $z$ is central in $\Aa$: $JzKv - vJzK = -Jzv$, so $e^{JzK}ve^{-JzK} = e^{-Jz}v$, and $e^{JzK}$ is invertible with inverse $e^{-JzK}$. At $z = -J\lambda$, $Jz = -J^{2}\lambda = \lambda$, so $e^{JzK} = e^{\lambda K} = R_{\lambda}$, the rotor flow of Theorem \claimref{singularity.modular-structure-completed-pair}(2), and $e^{-Jz} = e^{-\lambda}$, $e^{Jz} = e^{\lambda}$, in agreement with $R_{\lambda}vR_{\lambda}^{-1} = e^{-\lambda}v$ and $R_{\lambda}v^{*}R_{\lambda}^{-1} = e^{\lambda}v^{*}$ there. The script checks at two rational values $a = e^{\lambda/2}$ that $R = \tfrac12(a + a^{-1}) + \tfrac12(a - a^{-1})\sigma$ has inverse $\tfrac12(a + a^{-1}) - \tfrac12(a - a^{-1})\sigma$ and gives $RvR^{-1} = a^{-2}v$, $Rv^{*}R^{-1} = a^{2}v^{*}$ (the closed form of $e^{\lambda\sigma/2}$ from Theorem \claimref{std-algebra.bivector-exponential}), and that at $t = -J\lambda$, $-Jt = -\lambda$ and $JtK = \lambda K$.

\emph{(vi).} Write $x = x_{1} + x_{2}J$ with $x_{1}, x_{2} \in \Cl(1,1)$. Then $\br{x}_{0} = \br{x_{1}}_{0}$ and $\br{Jx}_{0} = -\br{x_{2}}_{0}$, so $T(x) = 2\br{x_{1}}_{0} + 2J\br{x_{2}}_{0}$, which is $T(Jx) = JT(x)$; $T(xy) = T(yx)$ follows from $\br{xy}_{0} = \br{yx}_{0}$ in any Clifford algebra and the centrality of $J$ in $\Aa$ ($\br{Jxy}_{0} = \br{yJx}_{0} = \br{Jyx}_{0}$). Under the isomorphism of (ii) $T$ is the matrix trace, since an $F$-linear functional on $\Mat_{2}(F)$ vanishing on commutators is a multiple of the trace and $T(1) = 2$; so $T(P_{\pm}) = 1$. $e^{-\beta K} = e^{-\beta/2}P_{+} + e^{\beta/2}P_{-}$ (the projectors are the spectral projectors of $\sigma$), so up to the positive factor $e^{-\beta/2}$, which cancels, $\varrho \propto P_{+} + yP_{-}$ with $y = e^{\beta} > 0$, $T(\varrho) \propto 1 + y$ and $\omega_{\beta}(v^{*}v) = \omega_{\beta}(P_{+}) = 1/(1 + y)$. Under the isomorphism with $\Mat_{2}$ over the field, in the Fock basis, $*$ is the conjugate transpose (it is a conjugate-linear antiautomorphism with $v \mapsto v^{*}$, so it agrees with the conjugate transpose on the generators $v$, $v^{*}$ and $J$) and $\varrho$ is diagonal with positive entries, so $\omega_{\beta}(x^{*}x) = T(\varrho x^{*}x)/T(\varrho)$ is the trace of the positive matrix $x\varrho x^{*}$ divided by a positive number: real and positive for $x \neq 0$. KMS: $\tau_{J\beta} = \operatorname{Ad}(e^{J\cdot J\beta K}) = \operatorname{Ad}(e^{-\beta K}) = \operatorname{Ad}(\varrho)$, and $T(\varrho\,x\,\varrho y\varrho^{-1}) = T(x\varrho y) = T(\varrho yx)$ by cyclicity, so $\omega_{\beta}(x\,\tau_{J\beta}(y)) = \omega_{\beta}(yx)$. At $-\beta$ the same computation gives $\omega_{\beta}(x\varrho^{-1}y\varrho) = T(\varrho x\varrho^{-1}y\varrho)/T(\varrho)$, which differs from $\omega_{\beta}(yx)$ for $x = v$, $y = v^{*}$ unless $y^{2} = 1$, that is $\beta = 0$. $\Aa$ is a matrix factor over the field $\mathrm{span}\{1, J\}$, and the flow is $\tau_{t} = \operatorname{Ad}(e^{JtK})$ with $J$ the unit of that field; Theorem \claimref{singularity.kms-uniqueness-factor} uses only the structure of a complex matrix factor, which $\Aa$ is with $J$ as its unit, so for $\beta \neq 0$ the state $\omega_{\beta}$ is the only $(\tau, \beta)$-KMS state. Finally $\ln\varrho = -\beta K - \ln T(e^{-\beta K})$, so $\varrho^{Jt} = e^{-J\beta tK}e^{-Jt\ln T(e^{-\beta K})}$, the second factor is central in $\Aa$, and $\operatorname{Ad}(\varrho^{Jt}) = \operatorname{Ad}(e^{J(-\beta t)K}) = \tau_{-\beta t}$, with $\tau_{-\beta t}(v) = e^{J\beta t}v$ by (iv). The script checks, at $y = 3$ and $y = 2/7$, that $T$ is $J$-valued on $1$ and $J$ and cyclic on all $64$ pairs of the real basis of $\Aa$; that $\omega_{\beta}(P_{+}) = 1/(1 + y)$; the KMS identity on all $64$ pairs, and its failure for $\operatorname{Ad}(\varrho^{-1})$; positivity of $\omega_{\beta}(x^{*}x)$ on ten elements; that the densities $r \in \Aa$ with $T(rx\varrho y\varrho^{-1}) = T(ryx)$ for all basis $x, y$ form a real space of dimension $2$, the $J$-line of $\varrho$; and, at a rational point $w = C + SJ$ of the unit circle of the field, that $\operatorname{Ad}(P_{+} + wP_{-})$ sends $v$ to $wv$ and $v^{*}$ to $\bar{w}v^{*}$, which is $\operatorname{Ad}(\varrho^{Jt})$ with $w = e^{J\beta t}$ up to a central factor. $\square$
